% Chapter Template

\section{Discussion} % Main chapter title

\label{Discussion} % Change X to a consecutive number; for referencing this chapter elsewhere, use \ref{ChapterX}

In this paper we have developed an extension of the agent based model of \cite{poledna_economic_2019} by replacing the constant consumption policy rule by the policy rule derived from the household optimization problem of \cite{kaplan_monetary_2018}. By doing this we introduce varying marginal propensities to consume which are observed in the data by \cite{fagereng_mpc_2021}. The agent based model is calibrated using micro and macro data from national accounts, sector accounts, input-output tables, government statistics, census data, business demography data, and household register data. Both the base agent based model and the adapted agent based model have been calibrated for the Dutch economy and have better long term forecasting performance to the benchmark models when forecasting the macroeconomic aggregates. The internal dynamics generated by the ABM are quite similar to their empirical counterparts well such as the cross-correlations of GDP and capital formation, household consumption, and operating surplus which implies that the model has both internal and external validity.\\

We have extended the ABM by introducing optimal consumption in household behavior. We have done this by using an adapted optimization strategy of \cite{maliar_deep_2021} which implements a neural network to solve the optimal consumption problem which the agents use to calculate their policy functions. This optimization strategy has two advantages. Firstly, as the method is not grid based the policy function is not restricted to a fixed state-space which was necessary due to the possible non stationary nature of the agent based model. Secondly, the optimization strategy is very fast. As the parameters which govern the expectations of the agent change every simulation period, the solution to the optimal consumption problem has to be recalculated every simulation period as well. Before switching to a neural network based method, we have tried grid based and collocation based optimization strategies and these strategies were simply too slow. For our forecasting analysis, we needed to solve the optimization problem 100,000 times and the grid based and collocation based methods would have taken weeks to calculate the new policy function each period for each simulation.\\

We found that including the optimal consumption model into the agent-based model improves forecasting performance for all forecasting exercises. The extended heuristic only improves forecasting performance in unconditional forecasting but performs worse with the unconditional ARX and conditional ARX forecasting exercise. \\

We also have analysed the internal dynamics of the agent based models. We found that the cross and auto-correlations of the variables generated by the agent based models are close to their empirical counterparts except for inflation. What we also found was that using ARX expectations greatly improved the likelihood of the model generated cross and auto-correlations to the data implying that this expectation mechanism might be a more realistic assumption. We also observed that the model generated cross-correlations of household consumption and wages have a timing error. This is most likely due to that wage setting follows production instead of being pro-cyclical and that household consumption follows wages closely.

The introduction of the adapted policy functions opens the way to calibrating the model on a 1:1 scale with household register data. This would mean that each agent would correspond to a Dutch household and we could exactly match their income and wealth levels and use this calibration strategy for more precise forecasting. Calibrating the model on household level would allow for estimating individual level direct and indirect policy effects by comparing actual data with predicted data with a policy experiment. 